%-----------------------------------
% Define document and include general packages
%-----------------------------------
% Tabellen- und Abbildungsverzeichnis stehen normalerweise nicht im
% Inhaltsverzeichnis. Gleiches gilt für das Abkürzungsverzeichnis (siehe unten).
% Manche Dozenten bemängeln das. Die Optionen 'listof=totoc,bibliography=totoc'
% geben das Tabellen- und Abbildungsverzeichnis im Inhaltsverzeichnis (toc=Table
% of Content) aus.
% Da es aber verschiedene Regelungen je nach Dozent geben kann, werden hier
% beide Varianten dargestellt.
\documentclass[12pt,oneside,titlepage,listof=totoc,bibliography=totoc]{scrartcl}
%\documentclass[12pt,oneside,titlepage]{scrartcl}

%-----------------------------------
% Dokumentensprache
%-----------------------------------
%\def\FOMEN{}% Auskommentieren um die Dokumentensprache auf englisch zu ändern
\newif\ifde
\newif\ifen

%-----------------------------------
% Meta informationen
%-----------------------------------
%-----------------------------------
% Meta Informationen zur Arbeit
%-----------------------------------

% Autor
\newcommand{\myAutor}{Julian Türner}

% Adresse
\newcommand{\myAdresse}{Heidestra\ss e 26 \\ \> \> \> 85356 Freising}

% Titel der Arbeit
\newcommand{\myTitel}{Stärken und Schwächen von Single Page Applikation Websites gegenüber traditionellen Webseiten}

% Betreuer
\newcommand{\myBetreuer}{Dr. rer. nat. Robert Heiniger}

% Lehrveranstaltung
\newcommand{\myLehrveranstaltung}{IT-Infrastruktur}

% Matrikelnummer
\newcommand{\myMatrikelNr}{581388}

% Ort
\newcommand{\myOrt}{München}

% Datum der Abgabe
\newcommand{\myAbgabeDatum}{\today}

% Semesterzahl
\newcommand{\mySemesterZahl}{4}

% Name der Hochschule
\newcommand{\myHochschulName}{FOM Hochschule für Oekonomie \& Management}

% Standort der Hochschule
\newcommand{\myHochschulStandort}{München}

% Studiengang
\newcommand{\myStudiengang}{Informatik}

% Art der Arbeit
\newcommand{\myThesisArt}{Literaturrecherche}

% Zu erlangender akademische Grad
\newcommand{\myAkademischerGrad}{Bachelor of Science (B.Sc.)}

% Firma
\newcommand{\myFirma}{Mustermann GmbH}


\ifdefined\FOMEN
%Englisch
\entrue
\usepackage[english]{babel}
\else
%Deutsch
\detrue
\usepackage[ngerman]{babel}
\fi

\newcommand{\langde}[1]{%
   \ifde\selectlanguage{ngerman}#1\fi}
\newcommand{\langen}[1]{%
   \ifen\selectlanguage{english}#1\fi}
\usepackage[utf8]{luainputenc}
\langde{\usepackage[babel,german=quotes]{csquotes}}
\langen{\usepackage[babel,english=british]{csquotes}}
\usepackage[T1]{fontenc}
\usepackage{fancyhdr}
\usepackage{fancybox}
\usepackage[a4paper, left=4cm, right=2cm, top=4cm, bottom=2cm]{geometry}
\usepackage{graphicx}
\usepackage{colortbl}
\usepackage[capposition=top]{floatrow}
\usepackage{array}
\usepackage{float}      %Positionierung von Abb. und Tabellen mit [H] erzwingen
\usepackage{footnote}
% Darstellung der Beschriftung von Tabellen und Abbildungen (Leitfaden S. 44)
% singlelinecheck=false: macht die Caption linksbündig (statt zentriert)
% labelfont auf fett: (Tabelle x.y:, Abbildung: x.y)
% font auf fett: eigentliche Bezeichnung der Abbildung oder Tabelle
% Fettschrift laut Leitfaden 2018 S. 45
\usepackage[singlelinecheck=false, labelfont=bf, font=bf]{caption}
\usepackage{caption}
\usepackage{enumitem}
\usepackage{amssymb}
\usepackage{mathptmx}
%\usepackage{minted} %Kann für schöneres Syntax Highlighting genutzt werden. ACHTUNG: Python muss installiert sein.
\usepackage[scaled=0.9]{helvet} % Behebt, zusammen mit Package courier, pixelige Überschriften. Ist, zusammen mit mathptx, dem times-Package vorzuziehen. Details: https://latex-kurs.de/fragen/schriftarten/Times_New_Roman.html
\usepackage{courier}
\usepackage{amsmath}
\usepackage[table]{xcolor}
\usepackage{marvosym}			% Verwendung von Symbolen, z.B. perfektes Eurozeichen

\renewcommand\familydefault{\sfdefault}
\usepackage{ragged2e}

% Mehrere Fussnoten nacheinander mit Komma separiert
\usepackage[hang,multiple]{footmisc}
\setlength{\footnotemargin}{1em}

% todo Aufgaben als Kommentare verfassen für verschiedene Editoren
\usepackage{todonotes}

% Verhindert, dass nur eine Zeile auf der nächsten Seite steht
\setlength{\marginparwidth}{2cm}
\usepackage[all]{nowidow}

%-----------------------------------
% Farbdefinitionen
%-----------------------------------
\definecolor{darkblack}{rgb}{0,0,0}
\definecolor{dunkelgrau}{rgb}{0.8,0.8,0.8}
\definecolor{hellgrau}{rgb}{0.0,0.7,0.99}
\definecolor{mauve}{rgb}{0.58,0,0.82}
\definecolor{dkgreen}{rgb}{0,0.6,0}

%-----------------------------------
% Pakete für Tabellen
%-----------------------------------
\usepackage{epstopdf}
\usepackage{nicefrac} % Brüche
\usepackage{multirow}
\usepackage{rotating} % vertikal schreiben
\usepackage{mdwlist}
\usepackage{tabularx}% für Breitenangabe

%-----------------------------------
% sauber formatierter Quelltext
%-----------------------------------
\usepackage{listings}
% JavaScript als Sprache definieren:
\lstdefinelanguage{JavaScript}{
	keywords={break, super, case, extends, switch, catch, finally, for, const, function, try, continue, if, typeof, debugger, var, default, in, void, delete, instanceof, while, do, new, with, else, return, yield, enum, let, await},
	keywordstyle=\color{blue}\bfseries,
	ndkeywords={class, export, boolean, throw, implements, import, this, interface, package, private, protected, public, static},
	ndkeywordstyle=\color{darkgray}\bfseries,
	identifierstyle=\color{black},
	sensitive=false,
	comment=[l]{//},
	morecomment=[s]{/*}{*/},
	commentstyle=\color{purple}\ttfamily,
	stringstyle=\color{red}\ttfamily,
	morestring=[b]',
	morestring=[b]"
}

\lstset{
	%language=JavaScript,
	numbers=left,
	numberstyle=\tiny,
	numbersep=5pt,
	breaklines=true,
	showstringspaces=false,
	frame=l ,
	xleftmargin=5pt,
	xrightmargin=5pt,
	basicstyle=\ttfamily\scriptsize,
	stepnumber=1,
	keywordstyle=\color{blue},          % keyword style
  	commentstyle=\color{dkgreen},       % comment style
  	stringstyle=\color{mauve}         % string literal style
}

%-----------------------------------
%Literaturverzeichnis Einstellungen
%-----------------------------------

% Biblatex

\usepackage{url}
\urlstyle{same}

%%%% Neuer Leitfaden (2018)
\usepackage[
backend=biber,
style=ext-authoryear-ibid, % Auskommentieren und nächste Zeile einkommentieren, falls "Ebd." (ebenda) nicht für sich-wiederholende Fussnoten genutzt werden soll.
%style=ext-authoryear,
maxcitenames=3,	% mindestens 3 Namen ausgeben bevor et. al. kommt
maxbibnames=999,
mergedate=false,
date=iso,
seconds=true, %werden nicht verwendet, so werden aber Warnungen unterdrückt.
urldate=iso,
innamebeforetitle,
dashed=false,
autocite=footnote,
doi=false,
useprefix=true, % 'von' im Namen beachten (beim Anzeigen)
mincrossrefs = 1
]{biblatex}%iso dateformat für YYYY-MM-DD

%weitere Anpassungen für BibLaTex
\input{skripte/modsBiblatex2018}

%%%%% Alter Leitfaden. Ggf. Einkommentieren und Bereich hierüber auskommentieren
%\usepackage[
%backend=biber,
%style=numeric,
%citestyle=authoryear,
%url=false,
%isbn=false,
%notetype=footonly,
%hyperref=false,
%sortlocale=de]{biblatex}

%weitere Anpassungen für BibLaTex
%\input{skripte/modsBiblatex}

%%%% Ende Alter Leitfaden

%Bib-Datei einbinden
\addbibresource{literatur/literatur.bib}

% Zeilenabstand im Literaturverzeichnis ist Einzeilig
% siehe Leitfaden S. 14
\AtBeginBibliography{\singlespacing}

%-----------------------------------
% Silbentrennung
%-----------------------------------
\usepackage{hyphsubst}
\HyphSubstIfExists{ngerman-x-latest}{%
\HyphSubstLet{ngerman}{ngerman-x-latest}}{}

%-----------------------------------
% Pfad fuer Abbildungen
%-----------------------------------
\graphicspath{{./}{./abbildungen/}}

%-----------------------------------
% Weitere Ebene einfügen
%-----------------------------------
\input{skripte/weitereEbene}

%-----------------------------------
% Paket für die Nutzung von Anhängen
%-----------------------------------
\usepackage{appendix}

%-----------------------------------
% Zeilenabstand 1,5-zeilig
%-----------------------------------
\usepackage{setspace}
\onehalfspacing

%-----------------------------------
% Absätze durch eine neue Zeile
%-----------------------------------
\setlength{\parindent}{0mm}
\setlength{\parskip}{0.8em plus 0.5em minus 0.3em}

\sloppy					%Abstände variieren
\pagestyle{headings}

%----------------------------------
% Präfix in das Abbildungs- und Tabellenverzeichnis aufnehmen, statt nur der Nummerierung (siehe Issue #206).
%----------------------------------
\KOMAoption{listof}{entryprefix} % Siehe KOMA-Script Doku v3.28 S.153
\BeforeStartingTOC[lof]{\renewcommand*\autodot{:}} % Für den Doppelpunkt hinter Präfix im Abbildungsverzeichnis
\BeforeStartingTOC[lot]{\renewcommand*\autodot{:}} % Für den Doppelpunkt hinter Präfix im Tabellenverzeichnis

%-----------------------------------
% Abkürzungsverzeichnis
%-----------------------------------
\usepackage[printonlyused]{acronym}

%-----------------------------------
% Symbolverzeichnis
%-----------------------------------
% Quelle: https://www.namsu.de/Extra/pakete/Listofsymbols.pdf
\usepackage[final]{listofsymbols}

%-----------------------------------
% Glossar
%-----------------------------------
\usepackage{glossaries}
\glstoctrue %Auskommentieren, damit das Glossar nicht im Inhaltsverzeichnis angezeigt wird.
\makenoidxglossaries
\newglossaryentry{server}{name={Server},description={Ein Server ist eine Software auf der eine Anwendungen ausgeführt wird}}
\newglossaryentry{browser}{name={Browser},description={Ein Browser ist eine Software auf einem Endgerät welches Quellcode der Webseiten für Nutzer darstellen kann}}
\newglossaryentry{schnittstelle}{name={Schnittstelle}, description={Eine Schnittstelle ist eine definierte Stelle, an der verschiedene Systeme miteinander kommunizieren können}}
\newglossaryentry{suchmaschine}{name={Suchmaschine}, description={Eine Suchmaschine durchsucht anhand von Schlagwörtern indizierte Inhalte}}
\newglossaryentry{controller}{name={Controller}, description={Ein Controller akzeptiert Anfragen und leitet diese Anfragen zur Bearbeitung weiter}}
\newglossaryentry{template}{name={Template}, description={Ein Template ist eine Vorlage für eine Form oder Struktur}}
\newglossaryentry{struktur}{name={Struktur}, description={Eine Struktur ist ein Aufbau aus mehreren Teilen}}
\newglossaryentry{overhead}{name={Overhead}, description={Ein Overhead ist ein Aufwand zuzüglich der primär Daten}}
\newglossaryentry{latenz}{name={Latenz}, description={Eine Latenz ist ein Zeitraum zwischen der Anfrage bis ein Ergebnis eintritt}}
\newglossaryentry{durchsatz}{name={Durchsatz}, description={Ein Durchsatz ist die Anzahl and verarbeiteten Daten in einem bestimmen Zeitraum}}
\newglossaryentry{keyword}{name={Keyword}, description={Ein Durchsatz ist die Anzahl and verarbeiteten Daten in einem bestimmen Zeitraum}}
\newglossaryentry{framework}{name={Framework}, description={Ein Framework stellt eine Reihe von Funktionen bereit, die verwendet werden können, um eine Anwendung zu entwickeln}}

%-----------------------------------
% PDF Meta Daten setzen
%-----------------------------------
\usepackage[hyperfootnotes=false]{hyperref} %hyperfootnotes=false deaktiviert die Verlinkung der Fußnote. Ansonsten inkompaibel zum Paket "footmisc"
% Behebt die falsche Darstellung der Lesezeichen in PDF-Dateien, welche eine Übersetzung besitzen
% siehe Issue 149
\makeatletter
\pdfstringdefDisableCommands{\let\selectlanguage\@gobble}
\makeatother

\hypersetup{
    pdfinfo={
        Title={\myTitel},
        Subject={\myStudiengang},
        Author={\myAutor},
        Build=1.1
    }
}

%-----------------------------------
% PlantUML
%-----------------------------------
%\usepackage{plantuml}

%-----------------------------------
% Umlaute in Code korrekt darstellen
% siehe auch: https://en.wikibooks.org/wiki/LaTeX/Source_Code_Listings
%-----------------------------------
\lstset{literate=
	{á}{{\'a}}1 {é}{{\'e}}1 {í}{{\'i}}1 {ó}{{\'o}}1 {ú}{{\'u}}1
	{Á}{{\'A}}1 {É}{{\'E}}1 {Í}{{\'I}}1 {Ó}{{\'O}}1 {Ú}{{\'U}}1
	{à}{{\`a}}1 {è}{{\`e}}1 {ì}{{\`i}}1 {ò}{{\`o}}1 {ù}{{\`u}}1
	{À}{{\`A}}1 {È}{{\'E}}1 {Ì}{{\`I}}1 {Ò}{{\`O}}1 {Ù}{{\`U}}1
	{ä}{{\"a}}1 {ë}{{\"e}}1 {ï}{{\"i}}1 {ö}{{\"o}}1 {ü}{{\"u}}1
	{Ä}{{\"A}}1 {Ë}{{\"E}}1 {Ï}{{\"I}}1 {Ö}{{\"O}}1 {Ü}{{\"U}}1
	{â}{{\^a}}1 {ê}{{\^e}}1 {î}{{\^i}}1 {ô}{{\^o}}1 {û}{{\^u}}1
	{Â}{{\^A}}1 {Ê}{{\^E}}1 {Î}{{\^I}}1 {Ô}{{\^O}}1 {Û}{{\^U}}1
	{œ}{{\oe}}1 {Œ}{{\OE}}1 {æ}{{\ae}}1 {Æ}{{\AE}}1 {ß}{{\ss}}1
	{ű}{{\H{u}}}1 {Ű}{{\H{U}}}1 {ő}{{\H{o}}}1 {Ő}{{\H{O}}}1
	{ç}{{\c c}}1 {Ç}{{\c C}}1 {ø}{{\o}}1 {å}{{\r a}}1 {Å}{{\r A}}1
	{€}{{\EUR}}1 {£}{{\pounds}}1 {„}{{\glqq{}}}1
}

%-----------------------------------
% Kopfbereich / Header definieren
%-----------------------------------
\pagestyle{fancy}
\fancyhf{}
% Seitenzahl oben, mittig, mit Strichen beidseits
% \fancyhead[C]{-\ \thepage\ -}

% Seitenzahl oben, mittig, entsprechend Leitfaden ohne Striche beidseits
\fancyhead[C]{\thepage}
%\fancyhead[L]{\leftmark}							% kein Footer vorhanden
% Waagerechte Linie unterhalb des Kopfbereiches anzeigen. Laut Leitfaden ist
% diese Linie nicht erforderlich. Ihre Breite kann daher auf 0pt gesetzt werden.
\renewcommand{\headrulewidth}{0.4pt}
%\renewcommand{\headrulewidth}{0pt}

%-----------------------------------
% Damit die hochgestellten Zahlen auch auf die Fußnote verlinkt sind (siehe Issue 169)
%-----------------------------------
\hypersetup{colorlinks=true, breaklinks=true, linkcolor=darkblack, citecolor=darkblack, menucolor=darkblack, urlcolor=darkblack, linktoc=all, bookmarksnumbered=false, pdfpagemode=UseOutlines, pdftoolbar=true}
\urlstyle{same}%gleiche Schriftart für den Link wie für den Text

%-----------------------------------
% Start the document here:
%-----------------------------------
\begin{document}

\pagenumbering{Roman}								% Seitennumerierung auf römisch umstellen
\newcolumntype{C}{>{\centering\arraybackslash}X}	% Neuer Tabellen-Spalten-Typ:
%Zentriert und umbrechbar

%-----------------------------------
% Textcommands
%-----------------------------------
\input{skripte/textcommands}

%-----------------------------------
% Titlepage
%-----------------------------------
\begin{titlepage}
	\newgeometry{left=2cm, right=2cm, top=2cm, bottom=2cm}
	\begin{center}
    \includegraphics[width=2.3cm]{abbildungen/fomLogo} \\
    \vspace{.5cm}
		\begin{Large}\textbf{\myHochschulName}\end{Large}\\
    \vspace{.5cm}
		\begin{Large}\langde{Hochschulzentrum}\langen{university location} \myHochschulStandort\end{Large}\\
		\vspace{2cm}
    \begin{Large}\textbf{\myThesisArt}\end{Large}\\
    \vspace{.5cm}
		% \langde{Berufsbegleitender Studiengang}
		% \langen{part-time degree program}\\
		% \mySemesterZahl. Semester\\
    \langde{im Studiengang}\langen{in the study course} \myStudiengang
		\vspace{1.7cm}

		% \langde{zur Erlangung des Grades eines}\langen{to obtain the degree of}\\
    \vspace{0.5cm}
	  	% \begin{Large}{\myAkademischerGrad}\end{Large}\\
		% Oder für Hausarbeiten:
		\textbf{im Rahmen der Lehrveranstaltung}\\
		\textbf{\myLehrveranstaltung}\\
		\vspace{1.8cm}
		\langde{über das Thema}
		\langen{on the subject}\\
    \vspace{0.5cm}
		\large{\textbf{\myTitel}}\\
		\vspace{2cm}
    \langde{von}\langen{by}\\
    \vspace{0.5cm}
    \begin{Large}{\myAutor}\end{Large}\\
	\end{center}
	\normalsize
	\vfill
    \begin{tabular}{ l l }
        \langde{Betreuer} % für Hausarbeiten
        %\langde{Erstgutachter} % für Bachelor- / Master-Thesis
        \langen{Advisor}: & \myBetreuer\\
        \langde{Matrikelnummer}
        \langen{Matriculation Number}: & \myMatrikelNr\\
        \langde{Abgabedatum}
        \langen{Submission}: & \myAbgabeDatum
    \\
    \end{tabular}
\end{titlepage}


%-----------------------------------
% Vorwort (optional; bei Verwendung beide Zeilen entkommentieren und unter Inhaltsverzeichnis setcounter entsprechend anpassen)
%-----------------------------------
%\input{kapitel/vorwort/vorwort}
%\newpage

%-----------------------------------
% Inhaltsverzeichnis
%-----------------------------------
% Um das Tabellen- und Abbbildungsverzeichnis zu de/aktivieren ganz oben in Documentclass schauen
\setcounter{page}{2}
\addtocontents{toc}{\protect\enlargethispage{-20mm}}% Die Zeile sorgt dafür, dass das Inhaltsverzeichnisseite auf die zweite Seite gestreckt wird und somit schick aussieht. Das sollte eigentlich automatisch funktionieren. Wer rausfindet wie, kann das gern ändern.
\setcounter{tocdepth}{4}
\tableofcontents
\newpage

%-----------------------------------
% Abbildungsverzeichnis
%-----------------------------------
% \listoffigures
% \newpage
%-----------------------------------
% Tabellenverzeichnis
%-----------------------------------
\listoftables
\newpage
%-----------------------------------
% Abkürzungsverzeichnis
%-----------------------------------
% Falls das Abkürzungsverzeichnis nicht im Inhaltsverzeichnis angezeigt werden soll
% dann folgende Zeile auskommentieren.
\addcontentsline{toc}{section}{\abbreHeadingName}

\section*{\langde{Abkürzungsverzeichnis}\langen{List of Abbreviations}}

\begin{acronym}[WYSIWYG]\itemsep0pt %der Parameter in Klammern sollte die längste Abkürzung sein. Damit wird der Abstand zwischen Abkürzung und Übersetzung festgelegt
  \acro{SPA}{Singe Page Application}
  \acro{HTML}{Hypertext Markup Language}
  \acro{CSS}{Cascading Style Sheets}
  \acro{JS}{JavaScript}
  \acro{API}{Application Programming Interface}
  \acro{SEO}{Search Engine Optimization}
\end{acronym}
\newpage

%-----------------------------------
% Symbolverzeichnis
%-----------------------------------
% In Overleaf führt der Einsatz des Symbolverzeichnisses zu einem Fehler, der aber ignoriert werdne kann
% Falls das Symbolverzeichnis nicht im Inhaltsverzeichnis angezeigt werden soll
% dann folgende Zeile auskommentieren.
% \phantomsection\addcontentsline{toc}{section}{\symheadingname}
% \input{skripte/symbolDef}
% \listofsymbols
% \newpage

%-----------------------------------
% Glossar
%-----------------------------------
\printnoidxglossaries
\newpage

%-----------------------------------
% Sperrvermerk
%-----------------------------------
%\input{kapitel/anhang/sperrvermerk}

%-----------------------------------
% Seitennummerierung auf arabisch und ab 1 beginnend umstellen
%-----------------------------------
\pagenumbering{arabic}
\setcounter{page}{1}

%-----------------------------------
% Kapitel / Inhalte
%-----------------------------------
% Die Kapitel werden über folgende Datei eingebunden
% Hinzugefügt aufgrund von Issue 167
%-----------------------------------
% Kapitel / Inhalte
%-----------------------------------
% Stellen Sie sicher, dass die Motivation und das Problem in einem klaren Zusammenhang zur Forschungsfrage stehen. Warum ist es wichtig, die Stärken und Schwächen von SPAs und traditionellen Webseiten zu untersuchen?

\section{Einleitung}
In der Einleitung wird das Thema der Arbeit eingeführt und die Motivation für die Arbeit dargelegt. 
Die Einleitung sollte den Leser dazu befähigen, die Arbeit in den Kontext der Forschungsliteratur einzuordnen.
Zudem sollte die Einleitung die Forschungsfrage und die Zielsetzung und einer Abgrenzung der Arbeit klar darlegen.
Darüber hinaus gibt die Einleitung den Aufbau der Arbeit wieder. 

\subsection{Hintergrund und Motivation}
% Einführung in das Thema der SPA-Websites und traditionellen Mehrseiten-Websites
% Bedeutung und Relevanz des Vergleichs der beiden Ansätze
% Zu Beginn des Internet, als das Internet noch in den Kinderschuhen steckte, bestanden Webseiten meist aus statischen \ac{HTML} Seiten.
% Die statischen Webseiten zeigten oft nur wie eine digitale Visitenkarte die Kontaktdaten, Produkte oder angebotene Services des Unternehmens an.
% Diese digitalen Visitenkarten zielten nicht darauf ab einen Profit mit und durch die Webseite zu erzielen. \cite[Seite 1]{Bly2018}
% Mittlerweile ist das Internet zu einem wichtigen Bestandteil des Lebens geworden und wird von vielen Menschen täglich genutzt.
% Durch die tägliche Nutzung wurde das Internet für viele zur primären Informationsquelle und einem wichtigen Kommunikationsmittel. \cite[Seite 1]{conf/pi/Sassenberg}
% Mit der transformation des Internet haben sich auch Webseiten verändert, welche sich mittlerweile in Anwendungen verwandelt haben und aktiv am Geschäft des Unternehmens beteiligt sind.
% Die neuen Anwendungen sollten für die Betreiber einen Profit erzielen, in dem sie das Branding des Unternehmens darstellen, Informationsquelle für die Nutzer sind vor allem aber eine starke Bindung zwischen dem Unternehmen und den Nutzern herstellen und aufrechterhalten. \cite[Seite 1]{Bly2018}
% Der Bedarf nach Webseiten und die Erwartungen an die Webseite scheint groß \cite[Seite 1]{Smith2022}, das belegt auch eine Umfrage der Community \Citeauthor{StackOverflow2023}, welche jährlich Softwareentwickler befragt welche Programmiersprachen und Technologien aktuell verwendet werden.
% In der Umfrage von \Citeauthor{StackOverflow2023} gaben  53\% der Befragten an, dass sie mit \ac{HTML} und \ac{CSS} arbeiten, was für die Entwicklung für Webseiten aller Art relevant ist. \cite[most-popular-technologies-programming-scripting-and-markup-languages]{StackOverflow2023}
% Die im Vergleich zu den digitalen Visitenkarten neuartigen Anwendungen werden als traditionelle Webseiten, Webanwendungen, Mehrseiten-Webseiten oder \ac{SPA} Webseiten bezeichnet, welche sich in ihrer Funktionsweise stark von den statischen Webseiten unterscheiden.
% Die traditionellen Webseiten bestehen aus mehreren \ac{HTML} Seiten, welche bei einem Aufruf auf einer Verknüpfung oder Schaltfläche eine komplett neue \ac{HTML} Seite laden.
% Die \ac{SPA} hingegen besteht aus einer einzigen \ac{HTML} Seite, welche bei einem Aufruf auf einer Verknüpfung oder Schaltfläche nur die Inhalte der Seite im \gls{browser} verändert und nicht zu einer neuen Seite navigiert.
% Durch diese Veränderung erhofft man sich Gewinne im Bereich der Performance, da die Webseite nicht bei jedem Aufruf eine neue Seite laden muss, sondern nur die Inhalte der Seite verändert.
% So wie mit jeder Technologie haben beide Ansätze verschiedene Umsetzungen und Implementierungen, was zu verschiedenen Meinungen und Ansichten von \ac{SPA} führt. \cite[Seite 4]{Flanagan2011}
% Das Ziel dieser Arbeit ist es, die Vor- und Nachteile der traditionellen Webseiten und der \ac{SPA} Webseiten herauszuarbeiten und zu vergleichen,
% um Entwicklern eine Entscheidungsgrundlage für die Entwicklung von Webseiten zu geben.


Mit der Entstehung des Internets bestanden Webseiten hauptsächlich aus statischen \ac{HTML} Seiten, die oft nur als digitale Visitenkarten dienten. 
Sie präsentierten grundlegende Informationen wie Kontaktdaten, Produkte oder Dienstleistungen eines Unternehmens, ohne den Fokus auf Profitabilität zu legen \cite{Bly2018}.
Doch mit der Zeit hat sich das Internet zu einem unverzichtbaren Bestandteil des täglichen Lebens entwickelt und dient vielen als primäre Informationsquelle und Kommunikationsmittel \cite{conf/pi/Sassenberg}.
Parallel dazu hat sich die Natur von Webseiten gewandelt. 
Sie sind von einfachen digitalen Visitenkarten zu komplexen Anwendungen herangewachsen, die aktiv zur Geschäftsstrategie eines Unternehmens beitragen. 
Diese Anwendungen zielen darauf ab, das Branding eines Unternehmens zu stärken, als Informationsquelle zu dienen und eine enge Bindung zwischen Unternehmen und Nutzern zu etablieren \cite{Bly2018}.
Eine Umfrage von \Citeauthor{StackOverflow2023} \cite{StackOverflow2023}, in der jährlich Softwareentwickler zu den von ihnen verwendeten Technologien befragt werden, zeigt, dass 53\% der Befragten mit \ac{HTML} und \ac{CSS} arbeiten.
Dies unterstreicht die Bedeutung der Webentwicklung in der heutigen Zeit.
Die modernen Webanwendungen können als traditionelle Webseiten, Webanwendungen, Mehrseiten-Webseiten oder \ac{SPA} kategorisiert werden. 
Während traditionelle Webseiten aus mehreren \ac{HTML}-Seiten bestehen, die bei einem Klick neu geladen werden, aktualisiert eine \ac{SPA} nur den Inhalt der aktuellen Seite. 
Dies verspricht Vorteile in der Performance. Doch wie bei jeder Technologie gibt es unterschiedliche Implementierungen und Meinungen zu \ac{SPA} \cite{Flanagan2011}.

Das Ziel dieser Arbeit ist es, die Stärken und Schwächen von traditionellen Webseiten im Vergleich zu \ac{SPA} zu analysieren und Entwicklern eine fundierte Entscheidungsgrundlage für die Webentwicklung zu bieten.


\subsection{Problemstellung und Zielsetzung der Literaturrecherche}
% Das Entwicklern einer Webseite oder einer Webanwendungen kann eine sehr komplexe Aufgabe sein, da es viele verschiedene Ansätze und Technologien gibt, welche sich selbst auch ständig weiterentwickeln.
% Durch die ständige Weiterentwicklung der Technologien und Ansätze ist es für Entwickler nicht immer einfach den Überblick zu behalten und die richtige Entscheidung zu treffen.
% Entwickler müssen sich mit den verschiedenen Ansätzen und Technologien auseinandersetzen und diese vergleichen, um die richtige Entscheidung für die Entwicklung der Webseite zu treffen.
% Dabei gilt es zu Verstehen was das Ziel der Webseite ist und welche Anforderungen die Webseite erfüllen muss um eine optimale Nutzererfahrung zu bieten.
% Die Thema dieser Arbeit ist es, die Vor- und Nachteile der traditionellen Webseiten und der \ac{SPA} Webseiten herauszuarbeiten und zu vergleichen.
% Die Problemstellung ist aus relevant für Entwickler, da die Vor- und Nachteile der beiden Ansätze einen Einfluss auf die Nutzererfahrung der Webseite haben.
% Die Forschungsfrage dieser Arbeit lautet: Welche Herausforderungen haben Entwickler bei \ac{SPA} Websites im Vergleich zu traditionellen Mehrseiten-Websites?

Die Entwicklung von Webseiten und Webanwendungen ist aufgrund der Vielzahl an Technologien und Ansätzen, die sich ständig weiterentwickeln, eine komplexe Aufgabe. 
Für Entwickler stellt sich die Herausforderung, stets auf dem neuesten Stand zu bleiben und die richtige Technologie oder den richtigen Ansatz für ihre Projekte auszuwählen. 
Insbesondere die Wahl zwischen traditionellen Mehrseiten-Webseiten und \ac{SPA} hat direkte Auswirkungen auf die Nutzererfahrung.

Das Hauptziel dieser Arbeit ist es, die Vor- und Nachteile von einer traditionellen Webseite im Vergleich zu einer \ac{SPA} zu analysieren. 
Dabei soll insbesondere untersucht werden, welche Herausforderungen Entwickler bei der Umsetzung einer \ac{SPA} im Vergleich zu traditionellen Mehrseiten-Webseiten begegnen.

\paragraph*{Forschungsfrage}
Welche Herausforderungen und Vorteile ergeben sich für Entwickler bei der Implementierung von einer \ac{SPA} im Vergleich zu traditionellen Mehrseiten-Websites?

\subsection{Abgrenzung}
Während diese Arbeit einen Einblick in die technischen und benutzerorientierten Aspekte von einer \ac{SPA} und traditionellen Webseiten bietet, werden Entwicklungskosten und spezifische Implementierungsdetails nicht im Detail behandelt.
Die Gründe für diese Abgrenzung sind die zahlreichen variablen Faktoren wie beispielsweise die Größe der Webseite, die Anzahl der Besucher und weitere spezifische Entwicklungsaspekte.

\subsection{Aufbau der Arbeit}
Die Arbeit ist in fünf Kapitel unterteilt, um den roten Faden zu wahren.
Im ersten Kapitel befindet sich die Einleitung, in welcher in das Thema der Arbeit eingeführt wird und die Motivation, Problemstellung sowie die Abgrenzung für die Arbeit dargelegt werden.
Das zweite Kapitel erläutert Grundlagen und Konzepte, welche für das Verständnis der Arbeit notwendig sind.
Das dritte Kapitel ist das Kapitel der Methodik, in der das Vorgehen nach \Citeauthor{conf/ecis/BrockeSNRPC09} beschreiben wird.
Im vierten Kapitel werden die Stärken und Schwächen von \ac{SPA} und traditionellen Webseiten aufgeführt, welche im Rahmen der Literaturrecherche gefunden wurden.
Neben der Aufführung wird die Forschungsfrage beantwortet und Implikationen der Ergebnisse für die Praxis aufgezeigt.
Die Arbeit endet mit dem Fazit, in dem die Arbeit reflektiert und ein Ausblick auf die Zukunft gegeben wird sowie die Ergebnisse der Arbeit zusammengefasst werden.

\input{kapitel/2_begriffe_grundlagen/begriffe_grundlagen.tex}
\section{Stärken und Schwächen von SPAs und traditionellen Webseiten}
\subsection{Schnittstellendesign und Benutzerinteraktion}
\subsection{Performance und Ladezeiten}
\subsection{Suchmaschinenoptimierung}

\input{kapitel/4_schlussbetrachtungen/schlussbetrachtungen.wlö}

%-----------------------------------
% Apendix / Anhang
%-----------------------------------
\newpage
\section*{\AppendixName} %Überschrift "Anhang", ohne Nummerierung
\addcontentsline{toc}{section}{\AppendixName} %Den Anhang ohne Nummer zum Inhaltsverzeichnis hinzufügen

\begin{appendices}
	% Nachfolgende Änderungen erfolgten aufgrund von Issue 163
	\makeatletter
	\renewcommand\@seccntformat[1]{\csname the#1\endcsname:\quad}
	\makeatother
	\addtocontents{toc}{\protect\setcounter{tocdepth}{0}} %
	\renewcommand{\thesection}{\AppendixName\ \arabic{section}}
	\renewcommand\thesubsection{\AppendixName\ \arabic{section}.\arabic{subsection}}
	% \section{Tabellen}
% 

\end{appendices}
\addtocontents{toc}{\protect\setcounter{tocdepth}{2}}

%-----------------------------------
% Literaturverzeichnis
%-----------------------------------
\newpage

% Die folgende Zeile trägt ALLE Werke aus literatur.bib in das
% Literaturverzeichnis ein, egal ob sie zietiert wurden oder nicht.
% Der Befehl ist also nur zum Test der Skripte sinnvoll und muss bei echten
% Arbeiten entfernt werden.
%\nocite{*}

%\addcontentsline{toc}{section}{Literatur}

% Die folgenden beiden Befehle würden ab dem Literaturverzeichnis wieder eine
% römische Seitennummerierung nutzen.
% Das ist nach dem Leitfaden nicht zu tun. Dort steht nur dass 'sämtliche
% Verzeichnisse VOR dem Textteil' römisch zu nummerieren sind. (vgl. S. 3)
%\pagenumbering{Roman} %Zähler wieder römisch ausgeben
%\setcounter{page}{4}  %Zähler manuell hochsetzen

% Ausgabe des Literaturverzeichnisses

% Keine Trennung der Werke im Literaturverzeichnis nach ihrer Art
% (Online/nicht-Online)
%\begin{RaggedRight}
%\printbibliography
%\end{RaggedRight}

% Alternative Darstellung, die laut Leitfaden genutzt werden sollte.
% Dazu die Zeilen auskommentieren und folgenden code verwenden:

% Literaturverzeichnis getrennt nach Nicht-Online-Werken und Online-Werken
% (Internetquellen).
% Die Option nottype=online nimmt alles, was kein Online-Werk ist.
% Die Option heading=bibintoc sorgt dafür, dass das Literaturverzeichnis im
% Inhaltsverzeichnis steht.
% Es ist übrigens auch möglich mehrere type- bzw. nottype-Optionen anzugeben, um
% noch weitere Arten von Zusammenfassungen eines Literaturverzeichnisse zu
% erzeugen.
% Beispiel: [type=book,type=article]
\printbibliography[nottype=online,heading=bibintoc,title={\langde{Literaturverzeichnis}\langen{Bibliography}}]

% neue Seite für Internetquellen-Verzeichnis
\newpage

% Laut Leitfaden 2018, S. 14, Fussnote 44 stehen die Internetquellen NICHT im
% Inhaltsverzeichnis, sondern gehören zum Literaturverzeichnis.
% Die Option heading=bibintoc würde die Internetquelle als eigenen Eintrag im
% Inhaltsverzeicnis anzeigen.
%\printbibliography[type=online,heading=bibintoc,title={\headingNameInternetSources}]
\printbibliography[type=online,heading=subbibliography,title={\headingNameInternetSources}]

\newpage
\pagenumbering{gobble} % Keine Seitenzahlen mehr

%-----------------------------------
% Ehrenwörtliche Erklärung
%-----------------------------------
\section*{%
	\langde{Ehrenwörtliche Erklärung}
	\langen{Declaration in lieu of oath}}
\langde{Hiermit versichere ich, dass die vorliegende Arbeit von mir selbstständig und ohne unerlaubte Hilfe angefertigt worden ist, insbesondere dass ich alle Stellen, die wörtlich oder annähernd wörtlich aus Veröffentlichungen entnommen sind, durch Zitate als solche gekennzeichnet habe. Ich versichere auch, dass die von mir eingereichte schriftliche Version mit der digitalen Version übereinstimmt. Weiterhin erkläre ich, dass die Arbeit in gleicher oder ähnlicher Form noch keiner Prüfungsbehörde/Prüfungsstelle vorgelegen hat. Ich erkläre mich damit einverstanden dass die Arbeit der Öffentlichkeit zugänglich gemacht wird. Ich erkläre mich damit einverstanden, dass die Digitalversion dieser Arbeit zwecks Plagiatsprüfung auf die Server externer Anbieter hochgeladen werden darf. Die Plagiatsprüfung stellt keine Zurverfügungstellung für die Öffentlichkeit dar.}


\par\medskip
\par\medskip

\vspace{5cm}

\begin{table}[H]
	\centering
	\begin{tabular*}{\textwidth}{c @{\extracolsep{\fill}} ccccc}
		\myOrt, \the\day.\the\month.\the\year
		&
		% Hinterlege deine eingescannte Unterschrift im Verzeichnis /abbildungen und nenne sie unterschrift.png
		% Bilder mit transparentem Hintergrund können teils zu Problemen führen
		\includegraphics[width=0.35\textwidth]{unterschrift}\vspace*{-0.35cm}
		\\
		\rule[0.5ex]{12em}{0.55pt} & \rule[0.5ex]{12em}{0.55pt} \\
		\langde{(Ort, Datum)}\langen{(Location, Date)} & \langde{(Eigenhändige Unterschrift)}\langen{(handwritten signature)}
		\\
	\end{tabular*} \\
\end{table}

\end{document}
